\documentclass[11pt]{article}
\usepackage[margin=1in]{geometry}
\usepackage{amsmath,amssymb,amsthm}
\usepackage{booktabs,array}
\usepackage{enumitem}
\usepackage[hidelinks]{hyperref}
\usepackage{url}
\setlength{\emergencystretch}{2em}
\setlist[itemize]{itemsep=2pt,leftmargin=*}
\setlist[enumerate]{itemsep=3pt,leftmargin=*}
\newtheorem{proposition}{Proposition}
\newtheorem{definition}{Definition}
\newcommand{\R}{\mathbb R}
\newcommand{\X}{\mathcal X}
\newcommand{\dd}{\mathrm d}
\newcommand{\cen}{\mathrm C}
\newcommand{\adv}{\mathrm{A}}
\newcommand{\argmax}{\operatorname*{arg\,max}}

\title{Fleet Coordination on Two Graphs:\\
\large Consolidated Findings and the Adversarial (Red-Team) Extension}
\author{}
\date{Working summary}

\begin{document}
\maketitle

\begin{abstract}
\noindent This is a one-stop summary of the model and results developed so far,
followed by the bridge to the adversarial question the reviewers asked us to
pursue. The central message is a single reusable idea: a distributed battery
fleet gains \emph{leverage over the grid and the price} only through three
couplings---endogenous price, shared network limits, and (for an adversary) a
detection budget---and that leverage turns on at a \emph{scale threshold}. The
same machinery that measures a coordination \emph{premium} for the owner
measures an adversary's \emph{capability}. The adversarial section is framed as
vulnerability mapping and detection on representative topologies, consistent
with a ``make the grid better, never worse'' posture; it is not an operational
attack guide.
\end{abstract}

\section{The model in one page}

\paragraph{Two coupled graphs.}
$G_1$ (physical) is a radial feeder with LinDistFlow flows; thermal and voltage
limits define the set $\mathcal A_t\subseteq\R^N$ of \emph{physically legal
joint injections} at time $t$. $G_2^i$ (strategy) is battery $i$'s
time-expanded state-of-charge graph; a path is a charge/discharge trajectory
$x_i=(c_i,g_i,e_i)$, with net injection $q_{it}=g_{it}-c_{it}$ and fleet
aggregate $Q_t=\sum_i q_{it}$. $G_1$ prunes joint moves; $G_2$ chooses among
what remains.

\paragraph{Price and payoff.}
Settlement price is aggregate-dependent,
\begin{equation}
p_t(Q_t)=f_t(L_t-Q_t),\qquad f_t\ \text{nondecreasing},\qquad
\text{affine case } p_t=a_t-\beta_t Q_t,\ \beta_t\ge 0 .
\label{eq:price}
\end{equation}
Realized fleet profit and the centralized problem are
\begin{equation}
\Pi(x)=\sum_t \Delta\,Q_t\,p_t(Q_t)-D(x),\qquad
x^{\cen}\in\argmax_{x\in\X}\Pi(x),
\label{eq:profit}
\end{equation}
with $D$ degradation cost and $\X$ the joint feasible set (local SoC/power
limits, household backup reserve, energy-neutral horizon, and
$q(t)\in\mathcal A_t$). Batteries are energy-limited (ours: $\approx$ 40\,kWh /
20\,kW, i.e.\ a two-hour asset), which caps sustained discharge.

\paragraph{Three regimes, one difference: the price each internalizes.}
\begin{center}
\begin{tabular}{@{}lll@{}}
\toprule
Regime & Marginal price term & Affine form\\
\midrule
Competitive price-taker & $p_t(Q_t)$ & $p_t$\\
Finite strategic battery $i$ & $p_t(Q_t)+q_{it}p_t'(Q_t)$ & $p_t-\beta_t q_{it}$\\
Central fleet & $p_t(Q_t)+Q_t p_t'(Q_t)$ & $p_t-\beta_t Q_t$\\
\bottomrule
\end{tabular}
\end{center}
The taker ignores its own impact; the planner internalizes the whole fleet's
impact $Q_t p_t'$. That single term is the self-cannibalization the fleet
avoids---and, later, the lever an adversary pulls.

\section{Key results (consolidated)}

\begin{enumerate}
\item \textbf{Competitive $=$ concave-potential maximizer.} The competitive
(price-taking / Wardrop) solution maximizes
$\Phi(x)=\sum_t\Delta\int_0^{Q_t}p_t(s)\dd s-\sum_i D_i$, which is concave; if
each $p_t$ is strictly decreasing the \emph{aggregate} $Q^{\mathrm W}$ is
unique (individual trajectories need not be). Correction to a common slip:
this potential's maximizer is the \emph{competitive} outcome, \emph{not} the
social optimum---conflating them forces the premium to $1$.

\item \textbf{Affine games have an exact potential.} Under $p_t=a_t-\beta_tQ_t$
the finite-player strategic game admits the exact potential
$\Psi(x)=\sum_t\Delta[a_tQ_t-\tfrac{\beta_t}{2}Q_t^2-\tfrac{\beta_t}{2}\sum_i
q_{it}^2]-\sum_iD_i$; global maximizers are (variational) Nash equilibria.

\item \textbf{Central dominance.} For any feasible baseline $x^{\mathrm B}$,
$\Pi(x^{\cen})\ge\Pi(x^{\mathrm B})$, so $\rho_{\mathrm B}=V_\cen/V_{\mathrm
B}\ge1$ when $V_{\mathrm B}>0$.

\item \textbf{Separability $\Rightarrow$ zero premium.} If prices are
exogenous, the feasible set is a product, and the objective is additive, then
central optimization equals the sum of independent optima---coordination adds
\emph{nothing}. A positive premium requires breaking one hypothesis:
\emph{shared network limits}, \emph{endogenous price}, or \emph{risk}.

\item \textbf{Affine impact bound.} With $K(x)=\tfrac{\Delta}{2}\sum_t\beta_t
Q_t^2$, the competitive-to-central gain obeys $0\le\Pi(x^{\cen})-\Pi(x^{\mathrm
W})\le K(x^{\mathrm W})$; the premium is small when impact slopes and dispatch
are small.

\item \textbf{No universal $4/3$ for profit.} A two-period affine example
($p_1=20+z,\ p_2=120-z$, capacity $C\uparrow50$) gives
$\rho_{\mathrm W}\to\infty$. The classical $4/3$ is a routing-\emph{cost}
bound and does not transfer to arbitrage profit.

\item \textbf{Scaling is an assumption, not a gift.} Individual price impact
vanishes as $N\to\infty$ (correction $\le\tfrac{\Delta\varepsilon_N}{2}\sum_t
\beta_tM_t\to0$), while the \emph{aggregate} correction $\beta_tQ_t$ can
persist. A deterministic mean-field limit is not implied by size alone;
common signals keep aggregate behavior correlated.
\end{enumerate}

\section{Coincident-peak / regulatory context}

The 4CP transmission regime is the institutional reason a single 15-minute
interval carries outsized weight, but for a \emph{residential} fleet it is
\emph{context, not a revenue line}: residential transmission is billed
volumetrically, so coincident-peak timing is monetized through the energy and
portfolio objectives already in the model, not a transmission credit. Cycling
with round-trip efficiency $<1$ slightly \emph{raises} total metered energy, so
the win is temporal energy-price spread, not avoided wires charges. The regime
is scheduled to change (Senate Bill 6; a 12CP / 30-minute proposal under PUCT
Project 58484, target Dec.~31 2026), so the setting is parameterized by an
interval count $m$ and length $\Delta_{\mathrm{CP}}$. ``Regulatory physics''
(nodal protocols) sits beside electrical physics: together they define what
efficient---or adversarial---actions look like.

\section{Numerical status (honest)}

\begin{center}
\begin{tabular}{@{}lrrrr@{}}
\toprule
Scenario & Units & Baseline (\$) & Central (\$) & $\Delta V$ (\$)\\
\midrule
Calm & 50{,}000 & 11{,}682 & 11{,}683 & 1\\
Calm & 150{,}000 & 34{,}876 & 34{,}884 & 8\\
Scarcity & 50{,}000 & 144{,}779 & 144{,}781 & 2\\
Scarcity & 150{,}000 & 323{,}061 & $\sim$325{,}000 & $\sim$1{,}939\\
\bottomrule
\end{tabular}
\end{center}

\begin{itemize}
\item At realistic ERCOT share the arbitrage coordination premium is
\textbf{$\approx1.00$}: energy arbitrage alone does not justify coordination.
\item A naive fleet piling into one spike \emph{overstates} its revenue by
$\sim$1.15--1.55$\times$; that is the cost of naivety, not a premium over a
smart baseline.
\item The premium-bearing regime (binding feeders + endogenous price at scale)
is not yet quantified---that simulation is the open item, and it is the same
simulation the adversarial question needs.
\end{itemize}

\section{Unifying thesis}

\begin{center}
\emph{Leverage over price and grid $=$ endogenous price ($\beta$) $+$ shared
network limits ($G_1$) $+$ (for an adversary) a detection budget, and it
switches on at a scale threshold.}
\end{center}
This one sentence organizes the owner's case (where coordination pays) and the
security case (where an adversary is dangerous). They are the same couplings
viewed with opposite objectives.

\section{The adversarial extension: same graphs, flipped objective}
\label{sec:adv}

We keep $G_1$, $G_2$, $\mathcal A_t$, and $\beta_t$ exactly as above. Only the
objective and one constraint change. This section is a red-team / vulnerability
formulation on representative or synthetic topologies; outputs are vulnerability
metrics, detection signatures, and mitigations---not tactics against specific
infrastructure.

\subsection{Adversary objectives}
A compromised fleet no longer maximizes $\Pi$. Two canonical objectives:
\begin{align}
\text{(Stealth cost inflation)}\quad
& J^{\adv}_{\text{cost}}(x)=\sum_t\Delta\,w_t\,\underbrace{p_t(Q_t)\,(L_t-Q_t)}_{\text{payment by load}},\\
\text{(Overt stress)}\quad
& J^{\adv}_{\text{stress}}(x)=\sum_t\sum_{e}\Big(\tfrac{|P_{et}(x)|}{K_{et}}\Big)^{+}
\ \ \text{or}\ \ \#\{(n,t):v_{nt}\notin[\underline V_n^2,\overline V_n^2]\}.
\end{align}
The stealth adversary maximizes system cost/volatility while remaining feasible
($x\in\X$, so protection never trips) and \emph{below detection}:
\begin{equation}
\big\|\,q_{it}-q_{it}^{\text{benign}}\,\big\|\le \varepsilon
\quad\text{per interval and per node (a detectability budget).}
\label{eq:stealth}
\end{equation}
The overt adversary drops \eqref{eq:stealth} and drives the joint action toward
the boundary of $\mathcal A_t$. Analysis identifies \emph{where} that boundary
is reachable; it is a map, not a recipe.

\subsection{Why the earlier results are already security results}
\begin{itemize}
\item \textbf{Separability $\to$ a no-harm zone.} With $\beta_t=0$ and no
binding feeder, the adversary's problem separates and its per-interval impact
is bounded by \eqref{eq:stealth}; correlated action buys it nothing. Harm, like
value, requires a coupling.
\item \textbf{Impact bound $\to$ capability bound.} $K(x)=\tfrac{\Delta}{2}
\sum_t\beta_tQ_t^2$ caps how far a given fleet can bend the price; it doubles
as the adversary's maximum cost-inflation lever at a given penetration.
\item \textbf{Self-cannibalization $\to$ the weapon.} The correlation the owner
diversifies away is exactly what the adversary maximizes to swing price or
stress a feeder. Coordinated-vs-independent harm is the mirror of the
coordination premium $\rho$.
\end{itemize}

\subsection{Three questions to answer (the deliverable)}
\begin{enumerate}
\item \textbf{Scale threshold.} Sweep fleet penetration $\phi$ (share of a
zone/feeder). Find $\phi^\ast$ at which the maximal price move or the first
$\mathcal A_t$ violation exceeds a materiality/detection threshold. This is the
reviewers' ``at what scale do market-making properties emerge,'' answered as a
security boundary.
\item \textbf{Correlation premium.} Quantify
$\text{harm}(\text{coordinated})/\text{harm}(\text{independent})$ vs.\ $\phi$:
how much extra damage comes purely from acting in concert.
\item \textbf{Stealth frontier.} Trace maximum harm as a function of the
detectability budget $\varepsilon$: the trade-off between staying invisible and
moving the system.
\end{enumerate}

\subsection{Defensive outputs}
\begin{itemize}
\item \textbf{Vulnerability map:} rank nodes/feeders by marginal harm
sensitivity---which $\mathcal A_t$ constraints sit closest to binding under an
adversarial $Q$, and which nodes have the largest local $\partial p/\partial
Q$. (Collapse the $\sim$20k-unit fleet into representative super-nodes; use
published representative grid topologies.)
\item \textbf{Detection signature:} the correlated-dispatch statistic that
separates adversarial from benign operation (benign diversifies; adversary
correlates)---a monitorable quantity.
\item \textbf{Mitigations:} per-feeder penetration caps, decorrelation
requirements, and telemetry-deviation bounds that map directly to
$\varepsilon$ in \eqref{eq:stealth}.
\end{itemize}

\subsection{Scoping note for the build}
A three-node toy is not representative; the interesting regime appears only at
scale. Minimal credible demonstration: a representative multi-node topology
with the fleet aggregated to super-nodes, real ERCOT price/load anchoring, and
one clean emergent result---a penetration threshold, a coordinated-vs-independent
divergence, or a detectable attack signature. UI is optional; the analysis is
the substance.

\end{document}